\scanpage{53}
\mspage{9}
% archive sheet number: 52
\uline{NB} Ayant (via un groupoïde équivalent à description purement combinatoire) ``reconstruit'' \unsure{disons} \deleted{\unsure{pour parler}} réalisé le groupoïde $S_*\mathscr T\unsure{\mathbb D|D}_{\mathrm{top}}^{\pm}$, et le foncteur
\diagram{\begin{tikzcd}[ampersand replacement=\&]
S_*\mathscr T\mathbb D_{\mathrm{top}}^{\pm}\arrow[r]\&(\mathrm{Gr.\ pond.\ rep.})
\end{tikzcd}}
\deleted{\unsure{alors}} on \unsure{devra|aura} \deleted{\unsure{par}} donc aussi \deleted{une description}\inserted{l'équivalent de} \deleted{de} la donnée d'un $(X,\Sigma,R)$, muni \uline{en plus} d'une partie $R'\subset R$ -- \deleted{\unsure{doit}} i.e. avoir la construction d'un groupoïde naturellement
\diagram{\begin{tikzcd}[ampersand replacement=\&]
{S_*\mathscr T\mathbb D_{\mathrm{top}}^{\pm}\ (\text{à sous ens. marqué }R'\text{ de sommets-repères})}\arrow[d]\\
S_*\mathscr T\mathbb D_{\mathrm{top}}^{\pm}
\end{tikzcd}}
et \deleted{d'un} il nous faudra donc \inserted{décrire} un foncteur ``gommage de $R'$''
\diagram{\begin{tikzcd}[ampersand replacement=\&,column sep=4em]
{S_*\mathscr T\mathbb D_{\mathrm{top}}^{\pm}(\text{avec }\deleted{\unsure{R}}R'\text{ marqué})}\arrow[r,"{\text{gommage de }R'}"]\arrow[r,shorten >=2em]\&S_*\mathscr T\mathbb D_{\mathrm{top}}^{\pm}
\end{tikzcd}}

% slanted
\margin{\uline{NB} Cette opération a l'air d'un \uline{affaiblissement} de structure -- mais cet affaiblissement, \unsure{pour|sur} une composante $C$ donnée de $\Sigma^+$, lorsqu'on ``gomme'' \uline{tous} les repères dans $C$ -- \unsure{sinon} c'est au contraire \unsure{plutôt dans la nature d'un} \unsure{renforcement} des structures, -- ou plutôt, des \unsure{morph.} \unsure{qui} \illegible{2} \unsure{dans $C$}, \unsure{si aucun des} $R'_C=R'\cap C$, \illegible{2}, \unsure{les} \unsure{catégories|opérations}\deleted{\unsure{ensembles}} de \unsure{gommage} \unsure{d'un côté} et celles du \unsure{surmarquage} (plus \unsure{bas}) \unsure{de l'autre} \unsure{sont inverses l'une de l'autre} \unsure{et sont des équivalences de groupoïdes}. Il faudrait \unsure{distinguer} les deux \unsure{cas}, pour distinguer les \uline{opérations de gommage} \uline{faibles} (\unsure{qui ne changent rien}) des \unsure{fortes}.}

\deleted{Cela s'exprime, sur le} donnant lieu à un diagramme comm. de foncteurs
\diagram{\begin{tikzcd}[ampersand replacement=\&,column sep=large,row sep=large]
{S_*\mathscr T\mathbb D_{\mathrm{top}}^{\pm}(\cdots R'\cdots)}\arrow[r,"{\text{gommage de }R'}"]\arrow[d,"{\text{oubli du marquage }R'}"']\&S_*\mathscr T\mathbb D_{\mathrm{top}}^{\pm}\arrow[dd]\&\\
S_*\mathscr T\mathbb D_{\mathrm{top}}^{\deleted{\unsure{+}}\pm}\arrow[dr]\&\&\\
\&(\mathrm{Gra}\deleted{\unsure{bi}}\mathrm{pond})\arrow[r]\&\substack{\text{graphes avec}\\\text{pondération}\\\text{par }\unsure{\gamma|g},\text{ mais}\\\text{non par }r,}
\end{tikzcd}}
alors qu'au niveau
